% FORMALIZACION_REUNIDA: CG-015/P09 del inventario REV05.
% Propietarios: sections/nucleo.tex (Cayley APP), libro/toro27.tex:123--256,
% PAQUETE_ARTICULO_I_INTEGRACION_ESTADO_CAMPOS_20260919/deltas/integracion/
% APPRouteWinding.lean (desplazamiento y cociclo entero).
% Esta insercion explicita el cambio de soporte y las fibras aritmeticas.
\section{Two-dimensional refinement of APP and carry cocycle}
\label{sec:refinamiento-app-27-9}

The realization of APP on a torus of nine positions per coordinate
and its expansion to twenty-seven positions are related through Euclidean
division. The expansion preserves two independent directions and assigns
an additional ternary digit to each coordinate. Thus each position
of the initial APP has nine lifts in the expanded support.
The two additional digits preserve the position within the fiber,
while the direction of transport and the regime of the TRIT reader
remain as distinct coordinates of the enriched state.

For $m\geq1$, we write
\[
 T_m=(\mathbb Z/m\mathbb Z)^2.
\]
In the calculations of this section the representatives
$0,\ldots,m-1$ are chosen. The nonadic presentation $1,\ldots,9$ prints the
zero class using the sign $9$. This convention allows the APP table
to be read and integer divisions to be performed with a single residue rule.

\subsection{Projection, fibers and ternary refinement coordinates}

The projection and the section of representatives are
\[
 p:T_{27}\longrightarrow T_9,\qquad
 p(x,y)=(x\bmod9,y\bmod9),
\]
\[
 s:T_9\longrightarrow T_{27},\qquad s(r,t)=(r,t),
 \quad 0\leq r,t<9.
\]
The projection is a surjective homomorphism of additive groups. Its
kernel is
\[
 \ker p=\{(9a,9b):a,b\in\mathbb F_3\}
        \simeq\mathbb F_3^2.
\]
The coordinate map
\begin{equation}
 \Psi:T_9\times\mathbb F_3^2\longrightarrow T_{27},\qquad
 \Psi((r,t),(a,b))=(r+9a,t+9b)
 \label{eq:app-refinamiento-coordenadas}
\end{equation}
uses representatives $0\leq a,b<3$.

\begin{proposition}[Unique decomposition of the expanded positions]
\label{prop:app-refinamiento-fibras}
The map $\Psi$ is bijective. Each fiber of $p$ contains nine
positions and the entire support satisfies
\[
 |T_{27}|=27^2=9^2\cdot3^2=729.
\]
The two coordinates of $\mathbb F_3^2$ are the two positional trits
that describe the refinement of the two APP coordinates.
\end{proposition}
\begin{proof}
For $0\leq x,y<27$, the Euclidean divisions
\[
 x=r+9a,\qquad y=t+9b,
 \quad 0\leq r,t<9,\quad0\leq a,b<3
\]
exist and are unique. These formulas construct the inverse of $\Psi$.
With $r,t$ fixed, the nine pairs $(a,b)$ are distinct and exhaust the
fiber. The census results from multiplying the $81$ residual positions
by their nine lifts.
\end{proof}

For example, the position $(5,5)$ has the fiber
\[
 \{(5+9a,5+9b):0\leq a,b<3\}.
\]
The row of this fiber with $b=1$ is
$(5,14),(14,14),(23,14)$. Each position preserves the same residual pair
$(5,5)$ and a different pair of quotients. Multiplicity nine expresses
two independent ternary choices. The topological and
temporal reading of TRIT additionally preserves its own operator domain.

\subsection{Additive composition and memory of modular crossing}

For $r,t\in\{0,\ldots,8\}$, define
\[
 r\oplus t=(r+t)\bmod9,\qquad
 c_9(r,t)=\left\lfloor\frac{r+t}{9}\right\rfloor.
\]
Transported by $\Psi$, addition in one coordinate of
$\mathbb Z/27\mathbb Z$ is
\begin{equation}
 (r,a)\boxplus(t,b)
   =\bigl(r\oplus t,\ a+b+c_9(r,t)\pmod3\bigr).
 \label{eq:app-refinamiento-adicion}
\end{equation}
In two coordinates, the same formula is applied separately to each axis.

\begin{proposition}[Cocycle of the positional extension]
\label{prop:app-refinamiento-cociclo}
The carry satisfies
\begin{equation}
 c_9(r,t)+c_9(r\oplus t,u)
 =c_9(t,u)+c_9(r,t\oplus u)
 =\left\lfloor\frac{r+t+u}{9}\right\rfloor.
 \label{eq:app-refinamiento-cociclo}
\end{equation}
Therefore, \eqref{eq:app-refinamiento-adicion} is associative and
exactly reproduces the additive group of the expanded support.
\end{proposition}
\begin{proof}
The equality $r+t=(r\oplus t)+9c_9(r,t)$ gives
\[
 \left\lfloor\frac{r+t+u}{9}\right\rfloor
 =c_9(r,t)+\left\lfloor\frac{(r\oplus t)+u}{9}\right\rfloor.
\]
The grouping $r+(t+u)$ provides the other side. Substituting
$x=r+9a$ and $y=t+9b$ into $x+y$ proves
\eqref{eq:app-refinamiento-adicion}; reducing its second
component modulo three corresponds to reducing $x+y$ modulo twenty-seven.
The cocycle identity directly proves the associativity of this
notation by residues and quotients.
\end{proof}

In particular, $8+1=0+9\cdot1$ shows the elementary crossing: the residue
returns to the zero class, printed as $9$ in the nonadic table, and the
first quotient increases by one. If only the residue
is preserved, the datum distinguishing $0$, $9$ and $18$ in the expanded
support is lost; the lifted state distinguishes these three positions.

The extension
\[
 0\longrightarrow\mathbb Z/3\mathbb Z
 \xrightarrow{\ a\mapsto9a\ }
 \mathbb Z/27\mathbb Z
 \longrightarrow\mathbb Z/9\mathbb Z\longrightarrow0
\]
is nonsplit as an extension of groups. Indeed, any
lift of the class $1$ is $1+9a$, of order twenty-seven. A
homomorphic section would have to send an element of order nine to one
whose order divides nine. The section of representatives instead has
the exact defect $9c_9$, which the preceding formula preserves.

\subsection{Compatibility of the multiplicative sheet}

The same decomposition preserves the second APP operation. For
one coordinate, we define
\[
 r\odot t=rt\bmod9,\qquad
 m_9(r,t)=\left\lfloor\frac{rt}{9}\right\rfloor.
\]
The lifted product is
\begin{equation}
 (r,a)\boxtimes(t,b)
 =\bigl(r\odot t,\ m_9(r,t)+rb+ta\pmod3\bigr).
 \label{eq:app-refinamiento-producto}
\end{equation}
\begin{proof}
The integer expansion provides
\[
 (r+9a)(t+9b)=rt+9(rb+ta)+81ab.
\]
The last term vanishes modulo twenty-seven. Substituting
$rt=(r\odot t)+9m_9(r,t)$ gives the formula. Uniqueness of Euclidean
division proves that its two components recover the residual
product and its refinement quotient.
\end{proof}
For example, $8\cdot8=64\equiv10\pmod{27}$ is represented by
$(1,1)$: the residue is $1$ and
$m_9(8,8)=7\equiv1\pmod3$. Thus sum and product preserve their
respective carries within the same expanded support.

\subsection{Directions and realization of transport}

Let $\mathcal D=\{N,S,E,O\}$ be the direction register and
$\delta_d\in\mathbb Z^2$ the unit displacement associated with $d$.
On the support with direction we define
\[
 S_m(z,d)=(z+\delta_d\pmod m,d).
\]
The projection explicitly preserves the direction,
$p_{\mathcal D}(z,d)=(p(z),d)$, and satisfies
\begin{equation}
 p_{\mathcal D}S_{27}=S_9p_{\mathcal D}.
 \label{eq:app-refinamiento-direccion}
\end{equation}
The quotient calculation uses, for an integer displacement $v$,
\[
 c_9(r;v)=\left\lfloor\frac{r+v}{9}\right\rfloor,
 \qquad r'=r+v-9c_9(r;v),\qquad
 a'=a+c_9(r;v)\pmod3.
\]
The formula includes negative displacements. For example,
$r=0,v=-1$ gives $r'=8,c_9=-1$; the reverse return restores the quotient.

\begin{proposition}[Conjugation of displacement in fiber coordinates]
The coordinate change $\Psi$ transforms $S_{27}$ into the
residual displacement $S_9$ accompanied by the update of the
two quotients just defined. This transformation is bijective
and preserves direction. Its powers preserve the same identity
for any finite length of transport.
\end{proposition}
\begin{proof}
On each axis, $x=r+9a$ is substituted into $x+v$ and Euclidean
division is performed. Direction remains fixed in the definition of $S_m$.
The inverse uses $-v$ and the cocycle, thereby recovering the initial
state. Composing the identity once per transition proves
the statement for all powers.
\end{proof}

The Hilbert realization preserves this decomposition. On the position
bases, the operator
\[
 \mathcal U\delta_{\Psi(z,a),d}
       =\delta_{z,d}\otimes\delta_a
\]
defines a surjective isometry
\[
 \ell^2(T_{27}\times\mathcal D)
 \simeq\ell^2(T_9\times\mathcal D)\otimes\ell^2(\mathbb F_3^2).
\]
The conjugated shift is the permutation operator with carry
described in the proposition. An additional internal factor $E$ remains
intact when tensoring with $I_E$. If a phase reader depends only on
the projection and a directional operator acts with the same matrix at
all positions, their compatibility squares also
commute. A reader using the quotient is realized precisely
in the preserved fiber factor. These conditions make explicit the
relation between the expanded support and the factors of the operator on
the torus, with the internal actions declared in its construction.

\subsection{Six ternary digits and distinct algebraic structures}

Each coordinate of $T_{27}$ has three base-three digits. Their joint
notation defines a bijection of sets
\[
 T_{27}\longleftrightarrow\mathbb F_3^6.
\]
This bijection publishes six positional digits. Addition on the torus
contains carries between digits and vector addition in
$\mathbb F_3^6$ is performed componentwise. For example, the
positional sum $2+1=3$ transforms $(2,0,0)$ into $(0,1,0)$, whereas
ternary vector addition would give $(0,0,0)$. The torus contains
elements of order twenty-seven; the ternary vector space has
exponent three. The two objects share a finite encoding,
with different typed operations. The fiber of ternary words,
the position support and the lines of a cubic surface
thus preserve their specific constructions and maps.
